\section{Marco teorico}

El analisis financiero es una lectura tecnica de los estados financieros que busca explicar la situacion de una empresa y apoyar decisiones economicas. El Marco Conceptual NIIF sostiene que la informacion financiera es util cuando sirve a los usuarios para tomar decisiones.

La NIC 1 destaca la importancia de la presentacion y comparabilidad de los estados financieros. Esta comparabilidad permite analizar cambios entre periodos y estudiar la estructura interna de los estados.

En el contexto peruano, el Plan Contable General Empresarial organiza cuentas contables que sirven como base para preparar estados financieros. Por ello, existe una relacion entre registro contable, clasificacion de cuentas, estados financieros y analisis.

Fuentes como Amat, Edwards, Hermanson y Maher, y la bibliografia contable registrada en el proyecto resaltan que el balance y el estado de resultados son instrumentos centrales para comprender la posicion financiera, el desempeno y la utilidad empresarial.
